\chapter{Osteoarthritis}
\index{Osteoarthritis}

\medskip
\noindent\textit{Educational information; seek professional advice for individual care.}

\section*{What it is}
Osteoarthritis is a condition in which tissues within and around a joint change over time. It commonly affects the knees, hips, hands, spine, or feet. Age, previous injury, joint shape, family history, repetitive loading, and body weight can influence risk; it is not simply unavoidable wear and tear.

\section*{Symptoms}
Pain with activity, short-lived stiffness after rest, swelling, reduced movement, tenderness, weakness, or a feeling that a joint is unstable are common. Symptoms and X-ray changes do not always match, and pain can also arise from other joint conditions.

\section*{Assessment}
Diagnosis is usually based on the history and examination. X-rays, blood tests, or joint-fluid testing may be used when the diagnosis is uncertain, symptoms are unusual, or infection, gout, inflammatory arthritis, or injury needs to be excluded.

\section*{Management}
Exercise adapted to ability, strengthening, joint-friendly activity, physical therapy, pacing, sleep support, and weight management where appropriate can improve pain and function. Medicines, topical treatments, injections, aids, or joint-replacement surgery may be considered according to the joint, risks, and response. Discuss medicine safety with a clinician, especially with kidney, stomach, heart, or liver disease.

\section*{When to seek urgent care}
Seek urgent help for a hot, very swollen joint with fever, sudden inability to bear weight after injury, new weakness or numbness, severe pain after a fall, or rapidly worsening illness. Arrange review when pain, stiffness, or loss of function interferes with daily activities.

\section*{Sources}
\begin{itemize}
\item National Institute of Arthritis and Musculoskeletal and Skin Diseases: \url{https://www.niams.nih.gov/health-topics/osteoarthritis}
\item NIAMS diagnosis and treatment: \url{https://www.niams.nih.gov/health-topics/osteoarthritis/diagnosis-treatment-and-steps-to-take}
\end{itemize}
